\section*{अध्याय \ref{ch:b1:membranes-transport} --- झिल्लियाँ और झिल्ली-अभिगमन}

\begin{solution}{exo:b1:membranes-transport:1}
फ़ॉस्फ़ोलिपिडों का एक द्विस्तर, पूँछें भीतर, कोलेस्टेरॉल के साथ, एक
\qty{5}{nm} चादर बनाता है। प्रोटीन उसके आरपार धँसे या किसी एक फलक से जुड़े
रहते हैं, अलग-अलग इकाइयों के रूप में। लिपिड और अधिकांश प्रोटीन तल में
पार्श्व विसरण करते हैं, जो एक द्विविमीय तरल है। दोनों पर्ण लिपिड-संघटन में
भिन्न हैं, और शर्करा-शृंखलाएँ केवल बाहरी फलक पर हैं।
\end{solution}

\begin{solution}{exo:b1:membranes-transport:2}
$\mathrm{O_2}$ (छोटा, अध्रुवी) $>$ इथेनॉल (छोटा, कमज़ोर ध्रुवी)
$>$ जल (छोटा, ध्रुवी) $\gg$ ग्लूकोज (बड़ा, ध्रुवी) $\gg$
$\mathrm{Na^+}$ (आवेशित, जलयोजित)। पारगम्यता ध्रुवता,
आकार और सबसे बढ़कर आवेश के साथ गिरती है, क्योंकि विलेय को
जलविरोधी गूदे में घुलना ही पड़ता है।
\end{solution}

\begin{solution}{exo:b1:membranes-transport:3}
सुक्रोस: $-2.48\times 0.1 = \qty{-0.25}{MPa}$। NaCl दो आयनों में
वियोजित होता है, यानी \qty{0.2}{osmol/L}: \qty{-0.50}{MPa}।
\end{solution}

\begin{solution}{exo:b1:membranes-transport:4}
$E_{Cl} = (61.5/(-1))\log(120/4) = -61.5\times 1.48 = \qty{-91}{mV}$।
\qty{-89}{mV} पर क्लोराइड साम्य से \qty{2}{mV} के भीतर है: वह
निष्क्रिय रूप से बँटा है।
\end{solution}

\begin{solution}{exo:b1:membranes-transport:5}
\omterm{def:b1:cell-unit-of-life:cell}{कोशिका} का पृष्ठ $\num{4.7e9}\times 10^{-10}\,\mathrm{m^2} =
\qty{0.47}{m^2}$; अनुपात $0.92/0.47 = 1.96 \approx 2$: लिपिड \omterm{def:b1:cell-unit-of-life:cell}{कोशिका} को
दो बार ढँकता है, इसलिए झिल्ली एक द्विस्तर है।
\end{solution}

\begin{solution}{exo:b1:membranes-transport:6}
\omterm{def:b1:cell-unit-of-life:cell}{कोशिका} का $\Psi = -0.9 + 0.5 = \qty{-0.4}{MPa}$; विलयन \qty{-0.6}{MPa}:
जल \omterm{def:b1:cell-unit-of-life:cell}{कोशिका} से निकलता है। साम्य तब जब $\Psi_{\text{कोशिका}} = -0.6$ हो:
$\Psi_p = -0.6 + 0.9 = \qty{0.3}{MPa}$; \omterm{def:b1:cell-unit-of-life:cell}{कोशिका} कुछ \omterm{def:b1:eukaryotic-cell:plantcell}{स्फीति} खोती है पर
स्फीत बनी रहती है।
\end{solution}

\begin{solution}{exo:b1:membranes-transport:7}
सांद्रता 10 से 20 तक दुगनी करने पर दर केवल \qty{15}{\%} चढ़ती है:
यानी संतृप्ति। $1/v$ को $1/c$ के सापेक्ष खींचने पर (या यह देखने पर कि
\qty{5}{mmol/L} पर $v = 1.0$ है और लगभग $1.67$ अनंत $c$ पर)
$V_{\max} \approx \qty{1.7}{\micro mol/min}$ मिलता है और अर्ध-अधिकतम लगभग
\qty{3}{mmol/L} पर: यानी किसी वाहक द्वारा सुगम विसरण।
\end{solution}

\begin{solution}{exo:b1:membranes-transport:8}
विश्राम पर झिल्ली की पारगम्यता पर खुले पोटैशियम \omterm{def:b1:membranes-transport:transporters}{चैनलों} का दबदबा रहता है;
पोटैशियम तब तक बाहर रिसता है जब तक भीतर उसे थामने लायक ऋणात्मकता न आ
जाए, यानी $E_K$ के पास। यदि सोडियम \omterm{def:b1:membranes-transport:transporters}{चैनल} खुल जाएँ, तो पारगम्यता पर
सोडियम का दबदबा होता और विभव $E_{Na}$ की ओर, यानी लगभग \qty{+60}{mV} की
ओर झूल जाता --- यही क्रिया-विभव है।
\end{solution}

\begin{solution}{exo:b1:membranes-transport:9}
भीतर आते $\mathrm{Na^+}$ के प्रति मोल: $RT\ln(12/145) + F(-0.060) = -6.4 - 5.8
= \qty{-12.2}{kJ}$; दो के लिए: \qty{-24.4}{kJ}। ग्लूकोज को चढ़ाई पर ले
जाने में $RT\ln(c_{\text{भीतर}}/c_{\text{बाहर}})$ लगता है; बराबरी से
$\ln(c_{\text{भीतर}}/c_{\text{बाहर}}) = 24\,400/2577 = 9.5$, यानी लगभग
\num{13000} का अनुपात।
\end{solution}

\begin{solution}{exo:b1:membranes-transport:10}
तरलता गिरती है (पूँछें जम जाती हैं, द्विस्तर जेल के पास आ जाता है);
ऑक्सीजन का विसरण धीमा होता है पर चलता रहता है; पंप, जो एक एंजाइम है,
लगभग रुक जाता है; रिसाव चलते रहते हैं, इसलिए घंटों में सोडियम घुसता है और
पोटैशियम निकलता है और प्रवणताएँ क्षीण हो जाती हैं; प्रवणताओं के मिटते ही
विभव शून्य की ओर गिरता है; और सोडियम के घुसने तथा प्रोटीनों के \omterm{def:b1:membranes-transport:osmosis}{परासरणी}
खिंचाव के अब संतुलित न रहने से जल घुसता है और \omterm{def:b1:cell-unit-of-life:cell}{कोशिका} फूल जाती है ---
ठंड में रखी \omterm{def:b1:cell-unit-of-life:cell}{कोशिकाएँ} और \omterm{def:b1:mammal-organization:organ}{अंग} ठीक इसी कारण फूलते हैं।
\end{solution}

\begin{solution}{exo:b1:membranes-transport:11}
\qty{37}{\celsius} पर मिलना पर \qty{4}{\celsius} पर न मिलना किसी
ताप-निर्भर प्रक्रिया का संकेत देता है --- यानी ऐसे तरल में विसरण जो ठंड
में ठोस हो जाता है। ATP के बिना न मिलना किसी सक्रिय, ऊर्जा-चालित पुनर्वितरण
(प्रेरक, पुटिका-चक्रण) का संकेत देता, विसरण का नहीं। सही परिणाम, कि
मिलने के लिए कोई ATP नहीं चाहिए, वही है जो तरल-मोज़ेक प्रतिरूप माँगता है:
प्रोटीन किसी द्विविमीय तरल में तापीय विसरण से चलते हैं, बिना \omterm{def:b1:cell-unit-of-life:cell}{कोशिका} के
कोई काम किए।
\end{solution}

\begin{solution}{exo:b1:membranes-transport:12}
विभव वह मान है जिस पर विद्युत बल सबसे अधिक पारगम्य आयन के विसरण को
संतुलित कर देता है, और उसे नेर्न्स्ट देता है; \omterm{def:b1:cell-unit-of-life:cell}{कोशिका} प्रवणताएँ तय करती है
(पंप से) और पारगम्यताएँ भी (अपने \omterm{def:b1:membranes-transport:transporters}{चैनलों} से), और विभव उनके पीछे चलता है।
\qty{-70}{mV} पर \qty{1}{\micro F/cm^2} की कोई झिल्ली
$\qty{7e-8}{C/cm^2}$ ढोती है, यानी प्रति वर्ग सेंटीमीटर लगभग
$\num{4e11}$ आयन --- किसी \qty{20}{\micro m} \omterm{def:b1:cell-unit-of-life:cell}{कोशिका} के लिए लगभग $10^5$
आयन, जबकि भीतर $10^{10}$ पोटैशियम आयन हैं: यानी आयनों का एक नगण्य अंश,
नगण्य दूरी चलकर, संधारित्र को आवेशित कर देता है।
प्रवणताएँ ही भंडार हैं; विभव उस पर की गई पढ़त है, और
पंप का सीधा योगदान (दो के बदले तीन आवेश बाहर) कुछ
मिलीवोल्ट है।
\end{solution}

\begin{solution}{pb:b1:membranes-transport:1}
\textbf{1.} $E_{Na} = 26.7\ln(145/12) = \qty{+66}{mV}$; $E_K =
26.7\ln(5/140) = \qty{-89}{mV}$।
\textbf{2.} $\mathrm{K^+}$ (\qty{29}{mV} दूर) $\mathrm{Na^+}$
(\qty{126}{mV} दूर) से अधिक पास है: झिल्ली $\mathrm{K^+}$ के प्रति कहीं
अधिक पारगम्य है।
\textbf{3.} भीतर आता $\mathrm{Na^+}$: $RT\ln(12/145) + F(-0.060) = -6.4 -
5.8 = \qty{-12.2}{kJ/mol}$। बाहर जाता $\mathrm{K^+}$: $RT\ln(5/140) -
F(-0.060) = -8.6 + 5.8 = \qty{-2.8}{kJ/mol}$।
\textbf{4.} प्रति चक्र लागत $= 3\times 12.2 + 2\times 2.8 = \qty{42.2}{kJ}$
जबकि प्रति ATP \qty{50}{kJ}: दक्षता \qty{84}{\%}।
\textbf{5.} पंप कोशिकाद्रव्यी सोडियम नीचा रखता है; शीर्ष फलक का \omterm{prop:b1:membranes-transport:secondary}{सहवाहक}
भीतर की ओर सोडियम-प्रवणता से ग्लूकोज को गुहा से भीतर खींचता है; और पंप को
दूसरे फलक पर होना ही चाहिए ताकि वह जो सोडियम बाहर करे वह रक्त में जाए, न
कि गुहा में वापस, जिससे गुहा का सोडियम बस चक्कर काटता रहता और शुद्ध
गति शून्य रह जाती।
\textbf{6.} प्रति मोल ग्लूकोज $2\times 12.2 = \qty{24.4}{kJ/mol}$।
\textbf{7.} $\Delta G = RT\ln(c_C/c_L)$ (ग्लूकोज अनावेशित है)।
\textbf{8.} $\ln(c_C/c_L) = 24\,400/2577 = 9.47$; $c_C/c_L =
\num{1.3e4}$।
\textbf{9.} सिद्धांततः \qty{1300}{mmol/L} तक: वह अनुपात कभी नहीं
पहुँचता क्योंकि ग्लूकोज आधारपार्श्व वाहक से निकल जाता है; व्यवहार में
\omterm{def:b1:cell-unit-of-life:cell}{कोशिकाद्रव्य} \qtyrange{5}{20}{mmol/L} के पास रहता है।
\textbf{10.} \omterm{def:b1:cell-unit-of-life:cell}{कोशिकाद्रव्य} (\qty{5}{mmol/L} से ऊपर) रक्त से अधिक सांद्र
है; वाहक ग्लूकोज को उस प्रवणता के नीचे रक्त में निष्क्रिय रूप से विसरित
होने देता है।
\textbf{11.} उलट दिया जाए तो \omterm{prop:b1:membranes-transport:secondary}{सहवाहक} ग्लूकोज को रक्त से \omterm{def:b1:cell-unit-of-life:cell}{कोशिका} में पंप
करता और वाहक उसे गुहा में रिसने देता: आँत ग्लूकोज स्रावित करती।
वाहकों की पक्षधरता ही अवशोषण की दिशा है।
\textbf{12.} \omterm{prop:b1:membranes-transport:secondary}{सहवाहक} सोडियम केवल ग्लूकोज के साथ लेता है; इसलिए गुहा का
ग्लूकोज उसके रास्ते सोडियम-अवशोषण चलाता है, तब भी जब हैजे का विष दूसरे
रास्ते रोक दे; क्लोराइड सोडियम के पीछे विद्युत रूप से चलता है और जल लवण
के पीछे \omterm{def:b1:membranes-transport:osmosis}{परासरण} से।
\textbf{13.} $\qty{4e-15}{mol/s}$ जितना $\mathrm{Na^+}$, यानी
प्रति सेकंड $\num{2.4e9}$ आयन; प्रति चक्र तीन पर, प्रति सेकंड $\num{8e8}$
चक्र।
\textbf{14.} प्रति सेकंड $\num{8e8}$ ATP $= \qty{1.3e-15}{mol/s}$; और
प्रति ग्लूकोज 30 ATP पर, \qty{4.4e-17}{mol/s} ग्लूकोज, यानी सोखे गए
ग्लूकोज का \qty{2.2}{\%}।
\textbf{15.} प्रति सेकंड $\num{1.2e9}$ ग्लूकोज, प्रति \omterm{prop:b1:membranes-transport:secondary}{सहवाहक} 50 पर:
$\num{2.4e7}$ \omterm{prop:b1:membranes-transport:secondary}{सहवाहक}; झिल्ली \qty{500}{\micro m^2}: प्रति वर्ग
माइक्रोमीटर \num{48000}, यानी प्रति \qty{20}{nm} वर्ग एक --- यानी
वाहकों से ठसाठस भरी कोई झिल्ली।
\textbf{16.} भीतर आते ऑस्मोल: $4\times\qty{2e-15}{mol/s} =
\qty{8e-15}{osmol/s}$; \qty{0.3}{osmol/L} पर, जल
$= \qty{2.7e-14}{L/s} = \qty{27}{\micro m^3/s}$। \omterm{def:b1:cell-unit-of-life:cell}{कोशिका} का आयतन
$25\times 5\times 5 = \qty{625}{\micro m^3}$: यानी \qty{23}{s} में दुगना।
जल उतनी ही तेज़ी से आधारपार्श्व झिल्ली के आरपार निकल जाता है जितनी तेज़ी
से घुसता है, और आँत जल इसी तरह सोखती है।
\textbf{17.} $\qty{2.7e-14}{L/s}$ जल यानी \qty{1.5e-12}{mol/s}; प्रति ग्लूकोज ($\qty{2e-15}{mol/s}$): लगभग 750 जल-अणु।
\textbf{18.} $S = 2(25\times 5 + 25\times 5 + 5\times 5) =
\qty{550}{\micro m^2} = \qty{5.5e-6}{cm^2}$; $C = \qty{5.5e-12}{F}$।
\textbf{19.} $Q = CV = 5.5\times 10^{-12}\times 0.060 =
\qty{3.3e-13}{C}$; $/\num{1.6e-19} = \num{2.1e6}$ आयन।
\textbf{20.} प्रति वर्ग माइक्रोमीटर $\num{2.1e6}/550 = 3800$ आयन, यानी झिल्ली के प्रति \qty{16}{nm} वर्ग एक अतिरिक्त आवेश।
\textbf{21.} \omterm{def:b1:cell-unit-of-life:cell}{कोशिका} में $\mathrm{K^+}$: $0.140\,\mathrm{mol/L}\times
\qty{6.25e-13}{L}\times\num{6e23} = \num{5.3e10}$ आयन; आवेश को उनमें से
\qty{0.004}{\%} चाहिए: यानी झिल्ली को आवेशित करने से सांद्रताएँ बदलती ही
नहीं।
\textbf{22.} प्रति चक्र शुद्ध आवेश $+e$ बाहर: धारा $\num{8e8}\times
\num{1.6e-19} = \qty{1.3e-10}{A}$; एक सेकंड में झिल्ली
\qty{1.3e-10}{C} कमा लेती, यानी $\Delta V = Q/C = \qty{24}{V}$। ऐसा
नहीं होता, क्योंकि बाहर पंप किया हर आवेश तुरंत उन आयनों से संतुलित हो
जाता है जो \omterm{def:b1:membranes-transport:transporters}{चैनलों} से वापस बहते हैं (मुख्यतः $\mathrm{K^+}$ कम रिसकर,
या $\mathrm{Na^+}$ फिर घुसकर): झिल्ली की चालकता पंप की धारा को लघुपथित
कर देती है, और पंप केवल मिलीवोल्ट जोड़ता है।
\textbf{23.} कोशिकाद्रव्यी सोडियम मिनटों में चढ़ जाता है; \omterm{prop:b1:membranes-transport:secondary}{सहवाहक} अपना
चालक बल खो देता है और ग्लूकोज का अवशोषण रुक जाता है; प्रवणताओं के चुकते
ही विभव शून्य की ओर क्षीण होता है; सोडियम तथा जल घुसते हैं और
\omterm{def:b1:cell-unit-of-life:cell}{कोशिका} फूल जाती है।
\textbf{24.} पोटैशियम-रिसाव रुक जाने पर सोडियम-रिसाव का दबदबा हो जाता
है: विभव $E_{Na}$ की ओर चलता है और कहीं कम ऋणात्मक हो जाता है
(विध्रुवण)।
\textbf{25.} लगभग \num{13000} का अनुपात (\omterm{def:b1:cell-unit-of-life:cell}{कोशिकाद्रव्य} पर गुहा), जो
\qty{-60}{mV} पर प्रति ग्लूकोज दो सोडियम आयनों से तय होता है; और यह
\omterm{def:b1:membranes-transport:transporters}{सोडियम--पोटैशियम पंप} (आधारपार्श्व) तथा सोडियम--ग्लूकोज \omterm{prop:b1:membranes-transport:secondary}{सहवाहक}
(शीर्ष) से संभव होता है।
\end{solution}
